\chapter*{R\'esum\'e}
\addcontentsline{toc}{chapter}{R\'esum\'e}
\markboth{R\'esum\'e}{R\'esum\'e}

Ce projet pr\'esente un prototype reproductible de d\'etection d'intrusions fond\'e
sur des observations de flux r\'eseau IoT. Le jeu de donn\'ees RT-IoT2022 est
contr\^ol\'e, d\'edupliqu\'e et transform\'e selon un sch\'ema versionn\'e. Une cascade
supervis\'ee associe un d\'etecteur binaire \`a un classifieur multiclasse de familles
d'attaque. Les candidats sont compar\'es sur des partitions distinctes
d'entra\^inement, de validation et de test, avec calibration et conservation des
artefacts promus.

Le prototype comprend une API FastAPI, une ingestion durable et idempotente,
une persistance transactionnelle, une diffusion temps r\'eel apr\`es validation et
une interface React destin\'ee \`a la supervision et \`a l'investigation. Une voie
Suricata distincte conserve les alertes issues de signatures sans les confondre
avec les pr\'edictions du mod\`ele. Le syst\`eme restitue aussi les versions des
artefacts, les preuves d'\'evaluation et des indicateurs de sant\'e du mod\`ele.

Sur le partage de test interne RT-IoT2022, la cascade atteint un macro-F1 de
0,9520 pour la graine~42 et une moyenne de 0,9686 sur trois graines. Ce r\'esultat
ne constitue toutefois pas une validation
sur un r\'eseau IoT r\'eel. L'absence de corpus NFStream natif \'etiquet\'e et
d'\'evaluation externe demeure la principale limite du travail.

\noindent\textbf{Mots-cl\'es :} Internet des objets, d\'etection d'intrusions,
flux r\'eseau, apprentissage automatique, RT-IoT2022, reproductibilit\'e.
